% Figure 6: Corralling for Catastrophic Failure Detection
% Adaptive Safety Mechanism in Multi-Expert LLM Routing

\section{Experiment: Corralling as Safety Mechanism}
\label{sec:corralling}

\subsection{Motivation: Fast Automatic Failover}

Production LLM systems face a critical challenge: \textbf{models can catastrophically fail with minimal warning}. API endpoints crash, providers silently update models with quality regressions, or traffic shifts to domains where warmup priors perform poorly. Manual detection and failover can take hours to days, during which users experience degraded service.

\textbf{The Question}: Can Corralling provide automatic, fast failover when a model catastrophically fails (e.g., API crashes, quality drops by 50\%+), without requiring human intervention?

\textbf{Connection to Prior Work}: This experiment complements our analysis of long-term adaptation (Figure~\ref{fig:corralling_semantic}), which focuses on adaptation timescales of 500--1,121 steps. This experiment tests \textbf{emergency response} in 3--50 steps. The warmup expert may contain semantic transfer priors (as in Figure~\ref{fig:corralling_semantic}), and this experiment validates that catastrophic failure detection works even when such priors are incorrect.

\subsection{The Corralling Algorithm: Exponential Reweighting with Safety}

Corralling~\cite{agarwal2017corralling} maintains a probability distribution over multiple expert policies and adaptively shifts weight based on observed losses. Unlike standard multi-armed bandits, it provides worst-case guarantees even when one expert is completely wrong.

\subsubsection{Exponential Weight Update with Exploration Floor}

\begin{equation}
p_{i,t+1} = (1-\gamma) \times \frac{p_{i,t} \cdot \exp(-\eta \cdot \hat{\ell}_{i,t})}{\sum_{j=1}^{K} p_{j,t} \cdot \exp(-\eta \cdot \hat{\ell}_{j,t})} + \frac{\gamma}{K}
\label{eq:corralling_mixed}
\end{equation}

where:
\begin{itemize}
    \item $p_{i,t}$: Probability of selecting expert $i$ at time $t$
    \item $\eta$: Learning rate (0.3 for fast catastrophic failure detection)
    \item $\hat{\ell}_{i,t}$: Importance-weighted loss estimate
    \item $\gamma$: Exploration floor (0.05, prevents expert death)
    \item $K$: Number of experts (2)
\end{itemize}

\textbf{Key modification}: The exploration floor $\gamma$ ensures every expert maintains $\geq \gamma/K$ probability, allowing recovery if environment changes.

\subsubsection{Importance-Weighted Loss Estimation}

\begin{equation}
\hat{\ell}_{i,t} = \begin{cases}
\frac{\ell_t}{p_{i,t}} & \text{if expert } i \text{ was selected} \\
0 & \text{otherwise}
\end{cases}
\label{eq:importance_weight}
\end{equation}

This estimator is unbiased: $\mathbb{E}[\hat{\ell}_{i,t}] = \ell_{i,t}$, enabling learning from bandit feedback.

\subsection{Experimental Setup: Three-Phase Catastrophic Failure}

\subsubsection{Design Rationale}

We test Corralling in its \textbf{correct operating regime}: catastrophic failures with large effect sizes (Cohen's $d > 1.5$), not subtle quality optimization ($d < 0.2$).

\textbf{Why this scenario?}
\begin{itemize}
    \item \textbf{Realistic}: Models crash, APIs timeout, providers update without notice
    \item \textbf{Urgent}: Requires fast detection (minutes/hours, not weeks)
    \item \textbf{Automatic}: Human-in-the-loop too slow for production
    \item \textbf{Right tool}: Corralling excels at large effects; offline A/B testing better for small effects
\end{itemize}

\subsubsection{Three-Phase Scenario}

We simulate a realistic production failure sequence:

\textbf{Phase 1 ($t=0-100$): Both Models Healthy}
\begin{itemize}
    \item Mixtral-8x7B: $\mu = 0.80$, $\sigma = 0.08$ (normal operation)
    \item GPT-4-Turbo: $\mu = 0.80$, $\sigma = 0.08$ (normal operation)
    \item Cohen's $d \approx 0$ (no difference)
    \item Expected: System maintains balanced weights
\end{itemize}

\textbf{Phase 2 ($t=100-300$): GPT-4 Catastrophically Fails}
\begin{itemize}
    \item Mixtral-8x7B: $\mu = 0.80$, $\sigma = 0.08$ (still healthy)
    \item GPT-4-Turbo: $\mu = 0.15$, $\sigma = 0.15$ (crashes, timeouts, errors)
    \item Cohen's $d \approx 5.0$ (catastrophic drop)
    \item Represents: API failures, severe quality regression
    \item Expected: Rapid decommissioning (3-50 steps)
\end{itemize}

\textbf{Phase 3 ($t=300-500$): GPT-4 Recovers}
\begin{itemize}
    \item Both models: $\mu = 0.80$, $\sigma = 0.08$ (provider fixed issue)
    \item Cohen's $d \approx 0$ (both healthy again)
    \item Tests: Can system detect recovery?
\end{itemize}

\subsubsection{Expert Policies}

\textbf{Warmup Expert (Stubborn):} Always selects GPT-4-Turbo (simulates rigid prior)

\textbf{Tabula Rasa Expert (Adaptive):} 95\% selects Mixtral-8x7B, 5\% explores (simulates online learner)

\textbf{Note}: We use deterministic mock experts for pedagogical clarity. Real LinUCB experts show more oscillations due to exploration (Supplementary Material).

\subsection{Results: Fast Catastrophic Failure Detection}

\subsubsection{Observed Dynamics (Figure~\ref{fig:corralling})}

Figure~\ref{fig:corralling} shows expert weight evolution over 500 routing decisions in the three-phase catastrophic failure scenario.

\begin{figure}[t]
\centering
\includegraphics[width=0.95\linewidth]{results/figure5_catastrophic_failure.pdf}
\caption{\textbf{Corralling for Catastrophic Failure Detection (Safety Regime).} 
\textbf{Phase 1 (t=0-100):} Both models healthy. System maintains balanced weights (40-60\%), demonstrating stability when evidence is ambiguous.
\textbf{Phase 2 (t=100-103):} GPT-4 catastrophically fails ($\mu: 0.80 \to 0.15$, Cohen's $d \approx 5.0$). System detects failure and decommissions warmup expert within 3 steps using moderate learning rate ($\eta=0.3$), achieving automatic failover without human intervention.
\textbf{Phase 3 (t=300-500):} GPT-4 recovers. System maintains conservative decommissioning with $\eta=0.3$. Recovery detection requires higher learning rates ($\eta=1.0$--$5.0$, as in Figure~\ref{fig:corralling_semantic}) or manual override.
\textbf{Two Operating Regimes}: This experiment demonstrates the \emph{safety regime} (3--50 steps, $\eta=0.3$), complementing the \emph{convergence regime} (300--1,121 steps, $\eta=5.0$, Figure~\ref{fig:corralling_semantic}).
\textbf{Robustness Validation}: Catastrophic failure detection ($t=3$--50) occurs $10\times$ faster than complete prior unlearning ($t \sim 300$--500, Figure~\ref{fig:corralling_semantic}), proving the system works \emph{even when warmup priors (including semantic transfer) are incorrect}.
\textbf{Key Result}: With large effect sizes ($d > 1.5$), Corralling achieves 100\% success with 3-50 step detection time, making it suitable for safety-critical production failover.}
\label{fig:corralling}
\end{figure}

\textbf{Phase 1 (t=0-100): Stability Under Normal Conditions}
- Both experts start at 50\% (uniform prior)
- Weights fluctuate 40-60\% due to sampling noise
- No premature decommissioning when both models work
- Validates: System doesn't collapse without evidence

\textbf{Phase 2 (t=100-103): Rapid Failure Detection}
- At $t=100$, GPT-4 quality drops catastrophically ($\mu: 0.80 \to 0.15$)
- Warmup expert immediately accumulates massive losses
- Exponential weight update with $\eta=0.3$ causes rapid decay
- \textbf{Decommissioning at $t=103$ (3 steps after failure)}
- Warmup weight drops below 10\% threshold
- Validates: Fast automatic failover (minutes in production)

\textbf{Phase 3 (t=300-500): Conservative Recovery (Learning Rate Dependent)}
- GPT-4 quality recovers ($\mu: 0.15 \to 0.80$)
- System maintains decommissioning (warmup weight $\approx 0\%$) with $\eta=0.3$
- \textbf{Recovery detection depends on learning rate regime:}
  \begin{itemize}
    \item $\eta=0.1$ (very conservative): No recovery detection, $>$800 steps needed
    \item $\eta=0.3$ (balanced, this experiment): Slow recovery, $>$500 steps
    \item $\eta=1.0$--$5.0$ (aggressive, Figure~\ref{fig:corralling_semantic}): Fast recovery detection, 50--200 steps
  \end{itemize}
- \textbf{Design trade-off}: Conservative $\eta=0.3$ prevents flapping (rapid on/off cycles) in safety-critical systems but requires manual override for recovery. For high-availability systems tolerant of transient failures, higher $\eta$ enables automatic recovery detection.
- Production recommendation: Use adaptive $\eta$—start at 0.3 for stability, increase to 1.0--2.0 after sustained failure to test recovery

\subsubsection{Performance Summary}

\textbf{Main Results (Catastrophic Scenario, $d \approx 5$):}
\begin{itemize}
    \item \textbf{Failure Detection}: 3 steps after catastrophic failure begins
    \item \textbf{Success Rate}: 100\% (across all tested seeds)
    \item \textbf{Sample Efficiency}: Only 500 total samples needed
    \item \textbf{Production Feasibility}: Hours to days (vs weeks for $d < 0.2$)
\end{itemize}

\subsection{Operating Regime Characterization}

To understand when Corralling provides value, we tested across multiple effect sizes (Supplementary Material):

\begin{table}[h]
\centering
\caption{Corralling Performance Across Effect Sizes}
\label{tab:corralling_regimes}
\begin{tabular}{lccc}
\toprule
\textbf{Scenario} & \textbf{Cohen's d} & \textbf{Detection Time} & \textbf{Use Case} \\
\midrule
Catastrophic failure & $> 1.5$ & 3-50 steps & ✅ API crashes \\
Severe degradation & $0.8-1.5$ & 100-300 steps & ✅ Version regression \\
Moderate mismatch & $0.3-0.8$ & 500-1000 steps & ⚠️ Domain shift \\
Subtle quality & $< 0.3$ & 2000+ steps & ❌ Use offline A/B \\
\bottomrule
\end{tabular}
\end{table}

\textbf{Key Insight}: Corralling's value proposition is \textbf{fast automatic failover for large effects}, not subtle quality optimization.

\subsection{Comparison to Alternatives}

\begin{table}[h]
\centering
\caption{Corralling vs Alternatives for Different Effect Sizes}
\label{tab:comparison}
\begin{tabular}{lccc}
\toprule
\textbf{Method} & \textbf{d > 1.5} & \textbf{d = 0.3-0.8} & \textbf{d < 0.3} \\
\midrule
\textbf{Corralling} & ✅ Hours & ⚠️ Days & ❌ Weeks \\
Offline A/B Test & ⚠️ 1 day & ✅ 3-7 days & ✅ 7-14 days \\
SPRT & ✅ Hours & ✅ 2-5 days & ⚠️ 10-20 days \\
\midrule
\textbf{Recommended} & Corralling/SPRT & Offline A/B & Offline A/B \\
\bottomrule
\end{tabular}
\end{table}

\textbf{Deployment Guideline}: Use Corralling when (a) catastrophic failures possible ($d > 1$), (b) high traffic (10k+ requests/day), or (c) cannot afford downtime for offline testing. Otherwise, use offline A/B testing.

\subsection{Supplementary Analysis}

\subsubsection{Realistic LMSYS Scenario (d = 0.12)}

To understand Corralling's limitations, we tested with realistic LMSYS reward distributions:
\begin{itemize}
    \item Mixtral: $\mu=0.823$, $\sigma=0.09$ (real data)
    \item GPT-4: $\mu=0.812$, $\sigma=0.10$ (real data)
    \item Cohen's $d \approx 0.12$ (tiny effect)
\end{itemize}

\textbf{Results (20 trials, 1,000 samples):}
\begin{itemize}
    \item Success rate: 25\% (5/20 trials)
    \item Mean detection: 526 $\pm$ 193 steps (when successful)
    \item Problem: Signal ($\Delta \mu = 0.011$) smaller than noise ($\sigma \approx 0.10$)
\end{itemize}

\textbf{With 10,000 samples}: Success rate increases to 100\%, mean detection 2,016 $\pm$ 1,992 steps.

\textbf{Conclusion}: Small effect sizes require 10x more samples (weeks/months in production), making offline A/B testing superior.

\subsubsection{Signal-to-Noise Analysis}

The realistic scenario fails due to four compounding factors:
\begin{enumerate}
    \item \textbf{Tiny effect}: $d = 0.12$ (100x smaller than catastrophic)
    \item \textbf{High overlap}: 46.7\% probability wrong ordering by chance
    \item \textbf{Insufficient samples}: Need 1,174 per expert, have only 500
    \item \textbf{Noise amplification}: Importance weighting amplifies noise by $\sqrt{1/p}$
\end{enumerate}

See Supplementary Material for detailed diagnostic analysis.

\subsection{Production Deployment Considerations}

\subsubsection{When to Use Corralling}

\textbf{Use when:}
\begin{itemize}
    \item High traffic (10,000+ requests/day) for fast convergence
    \item Large effect sizes ($d > 1.0$) such as API failures, crashes
    \item Safety-critical: need automatic failover without human intervention
    \item Stationary environment: task distribution relatively stable
\end{itemize}

\textbf{Don't use when:}
\begin{itemize}
    \item Low traffic (<1,000 requests/day) AND small effects ($d < 0.2$)
    \item Quality optimization: offline A/B testing faster and cheaper
    \item Non-stationary: frequent model updates or distribution shifts
    \item Can afford offline testing time (1-2 weeks)
\end{itemize}

\subsubsection{Deployment Constraints}

Five production constraints limit sample accumulation:
\begin{enumerate}
    \item \textbf{Time to convergence}: 10-100 days for low-traffic apps
    \item \textbf{Opportunity cost}: Lost quality during learning phase
    \item \textbf{Non-stationarity}: World changes over weeks/months
    \item \textbf{Context drift}: New topics, language evolution
    \item \textbf{Sample efficiency}: Offline testing often more efficient
\end{enumerate}

\textbf{Recommended approach}: Offline A/B test first ($d > 0.5$), then deploy with Corralling as safety net for catastrophic failures.

\subsection{Connection to Learning Rate Regimes}

\subsubsection{Two Operating Regimes Across Experiments}

Our experimental program (Figure~\ref{fig:corralling_semantic} and this figure) reveals two complementary operating regimes for Corralling with different learning rates:

\begin{table}[h]
\centering
\caption{Two Operating Regimes for Corralling-Based Systems}
\begin{tabular}{lccl}
\toprule
\textbf{Regime} & \textbf{Learning Rate} & \textbf{Timescale} & \textbf{Use Case} \\
\midrule
Safety (This Figure) & $\eta=0.3$--$1.0$ & 3--50 steps & Catastrophic failure detection \\
Convergence (Figure~\ref{fig:corralling_semantic}) & $\eta=2.0$--$5.0$ & 300--1,121 steps & Adapt beyond priors \\
\bottomrule
\end{tabular}
\end{table}

\paragraph{Unified Understanding:}
\begin{itemize}
    \item \textbf{Catastrophic failures (this experiment)} operate in the fastest regime. Detection time (3--50 steps) is much shorter than complete prior unlearning time ($\sim$300--500 steps, Figure~\ref{fig:corralling_semantic}). This validates robustness: \emph{even if warmup priors include incorrect semantic transfer}, catastrophic failures are detected before significant damage occurs.
    
    \item \textbf{Long-term adaptation (Figure~\ref{fig:corralling_semantic})} with aggressive learning ($\eta=5.0$) completely unlearns suboptimal priors, including semantic transfer. This demonstrates that catastrophic failure detection (fast, robust) and long-term convergence (slow, optimal) operate in different timescale regimes.
\end{itemize}

\subsubsection{Robustness to Incorrect Priors}

The warmup expert in this experiment may contain \textbf{semantic transfer} from related models (e.g., GPT-4o initialized from GPT-4-Turbo, as in Figure~\ref{fig:corralling_semantic}).

\textbf{Critical Validation}: For catastrophic failures (Cohen's $d > 1.5$):
\begin{itemize}
    \item Detection time: 3--50 steps (this experiment)
    \item Complete prior unlearning: $\sim$300--500 steps (Figure~\ref{fig:corralling_semantic})
    \item \textbf{Safety guarantee}: Failures detected $10\times$ faster than unlearning timeline
\end{itemize}

This demonstrates Corralling's robustness: catastrophic failure detection works \emph{even when warmup priors (including semantic transfer) are fundamentally incorrect}. The fast detection (3--50 steps) occurs well before the system would fully unlearn the prior ($\sim$300--500 steps with higher learning rates).

\subsection{Summary}

Corralling provides fast automatic failover for catastrophic model failures in production LLM systems. With large effect sizes ($d > 1.5$), it achieves 100\% success with 3-50 step detection time, making it suitable for safety-critical applications. However, for subtle quality differences ($d < 0.2$), offline A/B testing is superior due to sample efficiency and time-to-conclusion.

\textbf{Key Contributions}: 
\begin{enumerate}
    \item We characterize Corralling's operating regime, showing when online learning provides value (catastrophic failures) vs when offline testing is better (subtle quality optimization).
    \item We connect catastrophic failure detection (this figure, 3--50 steps) to the broader learning rate regime framework, including long-term convergence (Figure~\ref{fig:corralling_semantic}, $\eta=5.0$).
    \item We validate robustness: catastrophic failure detection works even when warmup priors (including potentially incorrect semantic transfer) are suboptimal, as detection occurs $10\times$ faster than the complete prior unlearning timeline.
\end{enumerate}

This deployment guidance is essential for practitioners choosing between online adaptation and offline testing strategies.

\subsection{Reproducibility}

Code available at: \texttt{experiments\_v1/06\_figure/}

\textbf{Main experiment}: \texttt{generate\_figure6\_main.py}
\begin{itemize}
    \item Configuration: $\eta=0.3$, $\gamma=0.05$, $N=500$ steps
    \item Three phases: Healthy (t=0-100) → Failure (t=100-300) → Recovery (t=300-500)
    \item Runtime: $\sim$3 seconds on MacBook Pro (M1)
\end{itemize}

\textbf{Supplementary experiments}:
\begin{itemize}
    \item Multi-seed validation: \texttt{supplementary/generate\_figure5\_multiseed.py}
    \item Realistic LMSYS: \texttt{supplementary/generate\_figure5\_realistic.py}
    \item Real LinUCB: \texttt{supplementary/generate\_figure5\_real\_linucb.py}
    \item Statistical power: \texttt{supplementary/test\_realistic\_10k\_samples.py}
    \item Diagnostic analysis: \texttt{supplementary/diagnostic\_realistic\_failure.py}
\end{itemize}

\bibliographystyle{ACM-Reference-Format}
\begin{thebibliography}{1}

\bibitem{agarwal2017corralling}
Alekh Agarwal, Haipeng Luo, Behnam Neyshabur, and Robert E. Schapire.
\newblock Corralling a band of bandit algorithms.
\newblock In \emph{Conference on Learning Theory (COLT)}, pages 12--38. PMLR, 2017.

\end{thebibliography}
